\documentclass[12pt]{article}
\usepackage[a4paper, margin=1in]{geometry}
\usepackage{hyperref}
\usepackage{enumitem}
\usepackage{array}

\title{Hybrid Filtering Based Book Recommender System Improvement Report}
\author{Darshan Nitin Barhate}
\date{July 2026}

\begin{document}
\maketitle

\section{Introduction}
This project is a Book Recommender System made using Python and Streamlit. It recommends books by using information from the Book-Crossing dataset. The main idea is to help users find books that match a title or author they like.

The improved version keeps the project simple, but makes it easier to run, easier to understand, and better documented.

\section{Project Objective}
The objective of this project is to build a simple web app that can:
\begin{itemize}[leftmargin=*]
    \item show popular books,
    \item recommend books based on title or author search,
    \item rank books using rating quality and rating count,
    \item show a small dataset summary,
    \item explain missing dataset files clearly.
\end{itemize}

\section{Technologies Used}
\begin{tabular}{|>{\raggedright\arraybackslash}p{4cm}|>{\raggedright\arraybackslash}p{9cm}|}
\hline
\textbf{Technology} & \textbf{Use in the Project} \\
\hline
Python & Main programming language used for the app. \\
\hline
Streamlit & Used to create the web interface. \\
\hline
Pandas & Used to read CSV files and calculate book ratings. \\
\hline
Book-Crossing Dataset & Provides books, ratings, and users data. \\
\hline
Jupyter Notebook & Used for study, testing, and recommender experiments. \\
\hline
\end{tabular}

\section{Dataset Files}
The project uses three CSV files:
\begin{itemize}[leftmargin=*]
    \item \texttt{Books.csv}: contains book title, author, publisher, year, and image URL.
    \item \texttt{Ratings.csv}: contains ratings given by users.
    \item \texttt{Users.csv}: contains user details.
\end{itemize}

These files should be placed inside an \texttt{archive} folder:
\begin{verbatim}
archive/
  Books.csv
  Ratings.csv
  Users.csv
\end{verbatim}

\section{Improvements Made}
\begin{enumerate}[leftmargin=*]
    \item \textbf{Clear missing file handling:} The app now shows a useful message if the dataset files are not present.
    \item \textbf{Cleaner data loading:} Duplicate ISBN values are removed and rating values are converted into numeric form.
    \item \textbf{Hybrid score:} Recommendations now use both average rating and rating count.
    \item \textbf{Useful sidebar controls:} The user can change minimum rating count and hide or show book covers.
    \item \textbf{Dataset summary page:} A new page shows total books, ratings, users, and top authors.
    \item \textbf{Better project setup:} A \texttt{requirements.txt} file was added for installing packages.
    \item \textbf{Better documentation:} A README and this LaTeX report were added.
    \item \textbf{Cleaner project files:} A \texttt{.gitignore} file was added to avoid uploading dataset and cache files.
\end{enumerate}

\section{How Recommendation Works}
First, the app calculates two values for every book:
\begin{itemize}[leftmargin=*]
    \item average rating,
    \item number of ratings.
\end{itemize}

When the user searches for a book title or author, the app finds matching books. Then it ranks those books using a simple hybrid score:
\begin{verbatim}
hybrid_score = average_rating * 0.7 + normalized_rating_count * 3
\end{verbatim}

This score gives importance to both book quality and popularity. A book with a high rating and many ratings will usually appear higher.

\section{Pages in the App}
\begin{itemize}[leftmargin=*]
    \item \textbf{Home:} Shows the project idea and featured popular books.
    \item \textbf{Recommend:} Lets the user search by title or author.
    \item \textbf{Popular Books:} Shows top-rated books with enough ratings.
    \item \textbf{Dataset Summary:} Shows dataset counts and top authors.
    \item \textbf{Settings:} Explains sidebar controls.
\end{itemize}

\section{How to Run}
Install the required packages:
\begin{verbatim}
pip install -r requirements.txt
\end{verbatim}

Run the Streamlit app:
\begin{verbatim}
streamlit run app.py
\end{verbatim}

\section{Testing and Validation}
The Python file was checked for syntax errors. The live Streamlit app requires the dataset files to be placed inside the \texttt{archive} folder.

\section{Conclusion}
The improved project is more complete and easier to understand. It now has better error handling, a simple hybrid ranking method, useful filters, a dataset summary page, and clear documentation. The project is suitable for a student-level machine learning or data science demonstration.

\end{document}
